\documentclass{../notatki}

\title{Teoria informacji kwantowej}

\begin{document}

\section{Wstęp}

Głównym źródłem zewnętrznym wiedzy jest
\textit{\href{https://www.preskill.caltech.edu/ph219/chap10_6A_2025.pdf}{wykład
Johna Preskilla}}.

Stan polaryzacji fotonu w systemie kwantowym, można opisać przy pomocy stanu:
$$
\ket{\psi} = \cos \theta \ket{\ |\ } + \sin \theta \ket{-}
$$
$$
\mathbb{P}(\ |\ ) = |\cos \theta|^2 \quad \mathbb{P}( - ) = |\sin \theta|^2
$$

Operator rzutowania $P$:
$$
P_0 = \ket{0}\bra{0} \quad P_1 = \ket{1}\bra{1}
$$
Ponieważ baza kubitu jest ortonormalna; $P_0 + P_1 = \ket{0}\bra{0} +
\ket{1}\bra{1} = 1$

$$
\mathbb{P}(0) = \ket{\psi} P_0 \bra{\psi}
$$

Stan układu po pomiarze to:
$$
\frac{P_0 \ket{\psi}}{\sqrt{\mathbb{P}(0)}}
$$

\end{document}
